\section{Diferenças Entre Virtualização e Emulação}

Algumas funcionalidades utilizadas neste trabalho - devido a limitações de hardware - serão emuladas em vez de virtualizadas.
A principal diferença entre esses dois conceitos está na forma em que as operações são executadas.

A emulação é realizada inteiramente por software, reproduzindo o comportamento de uma arquitetura ou dispositivo por meio da interpretação ou tradução
de instruções. Entretanto, devido à necessidade de traduzir instruções e simular componentes de hardware, a emulação tende a apresentar menor desempenho
em relação a virtualização.

Por outro lado, embora a virtualização também introduza uma camada de software, esta atua primariamente como uma interface de controle e gerenciamento.
Isso permite que a maior parte das instruções seja executada diretamente pelo processador físico. Contudo, essa abordagem exige suporte especializado
por parte do hardware subjacente, um requisito que nem sempre está disponível em todas as arquiteturas e máquinas.
